\section{Finite-sampling kernels for continuous Mellin synthesis}
\label{sec:tpc51-continuous-kernels}

We first separate a universal continuous-frame question from the
provenance-fixed Mellin direction used later.  The former has a sharp
negative answer for a reason that is independent of arithmetic: at each
finite value of \(X\), the physical output space is finite dimensional,
whereas a freely weighted continuous dictionary has infinitely many
coefficient directions.

\subsection{A smooth kernel in every finite-dimensional output}

Let \(\Theta\subset\R^d\) be open, let \(\cH\) be a complex Hilbert
space of finite dimension \(N\), and let

\[
 Z:\Theta\longrightarrow\cH
\]

be continuous.  For \(g\in C_c^\infty(\Theta)\), define
\begin{align}
 \cS_Zg&:=\int_\Theta g(\tau)Z(\tau)\,d\tau,
 \label{eq:tpc51-continuous-synthesis}\\
 \cE_Z(g)&:=\int_\Theta |g(\tau)|\,\|Z(\tau)\|_\cH\,d\tau.
 \label{eq:tpc51-continuous-envelope}
\end{align}
The first integral is a Bochner integral.  Both quantities are finite
because \(g\) has compact support.

\begin{theorem}[Universal smooth-kernel obstruction]
\label{thm:tpc51-smooth-kernel}
Suppose \(Z\) is not identically zero.  Then there exists a nonzero
\(g\in C_c^\infty(\Theta)\) such that
\begin{equation}
 \cS_Zg=0,
 \qquad
 \cE_Z(g)>0.
 \label{eq:tpc51-smooth-kernel}
\end{equation}
Consequently, there is no finite constant \(C\) for which
\begin{equation}
 \cE_Z(g)\le C\|\cS_Zg\|_\cH
 \label{eq:tpc51-false-universal-frame}
\end{equation}
holds for every \(g\in C_c^\infty(\Theta)\).
\end{theorem}

\begin{proof}
Choose \(\tau_0\in\Theta\) with \(Z(\tau_0)\ne0\).  By continuity,
there is an open ball \(U\Subset\Theta\) and a number \(m>0\) such that
\(\|Z(\tau)\|_\cH\ge m\) on \(U\).  Choose \(N+1\) pairwise disjoint
balls \(U_k\Subset U\) and nonzero functions
\(\rho_k\in C_c^\infty(U_k)\).  Put
\[
 w_k:=\int_\Theta \rho_k(\tau)Z(\tau)\,d\tau\in\cH.
\]
The \(N+1\) vectors \(w_k\) are linearly dependent, so there are
coefficients \(c_k\), not all zero, such that
\(\sum_kc_kw_k=0\).  The function
\(g=\sum_kc_k\rho_k\) is nonzero because the supports of the
\(\rho_k\) are disjoint, and \(\cS_Zg=0\).  Moreover,
\[
 \cE_Z(g)
 \ge m\sum_k |c_k|\int_{U_k}|\rho_k(\tau)|\,d\tau>0.
\]
This proves \eqref{eq:tpc51-smooth-kernel};
\eqref{eq:tpc51-false-universal-frame} is then impossible.
\end{proof}

The theorem remains true when \(N=N(X)\) grows: it is applied
separately at each finite \(X\).  It does not use an artificial
refactorization of a fixed physical tensor.  It states that the
coefficient space in \eqref{eq:tpc51-continuous-synthesis} cannot be
the whole of \(C_c^\infty(\Theta)\) if a lower frame is desired.

\subsection{An explicit kernel for logarithmic phase samples}

The preceding dimension argument can be made completely explicit for
the scalar Mellin phase dictionary.  For \(h\in C_c^\infty(\R)\), use
the positive-sign Fourier convention
\begin{equation}
 \cF_+h(y):=\int_\R h(t)\e^{ity}\,dt.
 \label{eq:tpc51-positive-Fourier}
\end{equation}
Integration by parts gives the sign-sensitive identity
\begin{equation}
 \cF_+\bigl((\partial_t+iy_0)h\bigr)(y)
 =i(y_0-y)\cF_+h(y).
 \label{eq:tpc51-ibp-sign}
\end{equation}

\begin{proposition}[Explicit phase-sampling kernel]
\label{prop:tpc51-explicit-phase-kernel}
Let \(y_1,\ldots,y_N\in\R\), let
\(a_1,\ldots,a_N\in\C\setminus\{0\}\),
and put
\begin{equation}
 Z_{\bm y}(t)
 :=\bigl(a_ne^{ity_n}\bigr)_{n=1}^N\in\C^N.
 \label{eq:tpc51-phase-dictionary}
\end{equation}
For every nonzero \(\rho\in C_c^\infty(\R)\), the function
\begin{equation}
 g_{\bm y,\rho}(t)
 :=\prod_{n=1}^N(\partial_t+iy_n)\rho(t)
 \label{eq:tpc51-explicit-phase-kernel-function}
\end{equation}
is nonzero and satisfies
\begin{equation}
 \int_\R g_{\bm y,\rho}(t)Z_{\bm y}(t)\,dt=0.
 \label{eq:tpc51-explicit-phase-kernel-output}
\end{equation}
\end{proposition}

\begin{proof}
Iterating \eqref{eq:tpc51-ibp-sign} gives
\begin{equation}
 \cF_+g_{\bm y,\rho}(y)
 =i^N\prod_{n=1}^N(y_n-y)\cF_+\rho(y).
 \label{eq:tpc51-explicit-kernel-multiplier}
\end{equation}
The \(k\)-th coordinate of
\eqref{eq:tpc51-explicit-phase-kernel-output} is
\(a_k\cF_+g_{\bm y,\rho}(y_k)=0\).

It remains to check that \(g_{\bm y,\rho}\ne0\).  If it vanished,
\eqref{eq:tpc51-explicit-kernel-multiplier} would imply that
\(\cF_+\rho(y)=0\) for every real
\(y\notin\{y_1,\ldots,y_N\}\).
The Fourier transform of a compactly supported smooth function extends
to an entire function.  It would therefore vanish identically, forcing
\(\rho=0\), a contradiction.
\end{proof}

There is an equivalent physical-shape formulation.  If the logarithmic
sample points are \(y_n\), choose a nonzero
\(F_X\in C_c^\infty(\R)\) of the form
\begin{equation}
 F_X(s)=\chi(s)\prod_{n=1}^N(s-y_n).
 \label{eq:tpc51-sample-zero-shape}
\end{equation}
Its Fourier transform is Schwartz, while Fourier inversion gives
\(F_X(y_n)=0\) at every sample.  This example depends on the full sample
set, and its derivative seminorms need not be uniform in \(X\).  It
therefore rules out a theorem uniform over arbitrary smooth shapes; it
does not contradict a theorem for the fixed provenance shapes.

\subsection{Near kernels and phase aliases}

For a finite complex Borel measure \(\nu\) on \(\Theta\), extend the
notation by
\begin{equation}
 \cS_Z\nu:=\int_\Theta Z(\tau)\,d\nu(\tau),
 \qquad
 \cE_Z(\nu):=\int_\Theta\|Z(\tau)\|_\cH\,d|\nu|(\tau).
 \label{eq:tpc51-measure-synthesis}
\end{equation}

\begin{proposition}[Neighboring-atom condition number]
\label{prop:tpc51-neighboring-atoms}
\label{prop:tpc51-neighbouring-atoms}
Let \(I\subset\R\) be an interval on which \(Z:I\to\cH\) is
continuously differentiable, and suppose
\begin{equation}
 \inf_{t\in I}\|Z(t)\|_\cH\ge m>0,
 \qquad
 \sup_{t\in I}\|Z'(t)\|_\cH\le L
 \quad\text{for some }L>0.
 \label{eq:tpc51-atom-regularity}
\end{equation}
If \(t,t+h\in I\) and \(h\ne0\), then for
\(\nu_{t,h}=\delta_t-\delta_{t+h}\),
\begin{equation}
 \frac{\cE_Z(\nu_{t,h})}
      {\|\cS_Z\nu_{t,h}\|_\cH}
 \ge\frac{2m}{L|h|},
 \label{eq:tpc51-neighbor-condition}
\end{equation}
with the ratio interpreted as \(+\infty\) when the denominator is zero.
\end{proposition}

\begin{proof}
The two atoms are distinct, so
\(\cE_Z(\nu_{t,h})=\|Z(t)\|+\|Z(t+h)\|\ge2m\).
The fundamental theorem of calculus gives
\[
 \|\cS_Z\nu_{t,h}\|
 =\|Z(t)-Z(t+h)\|
 \le L|h|.
\]
\end{proof}

Thus continuity of a dictionary is an obstruction, rather than a
source of a uniform signed lower frame, when its coefficient measure is
free.  Narrow smooth bumps centered at \(t\) and \(t+h\) give the same
conclusion up to an arbitrarily small relative error.

For completeness, finite phase sampling also has long-range aliases.

\begin{proposition}[Simultaneous phase alias]
\label{prop:tpc51-simultaneous-alias}
For the dictionary \eqref{eq:tpc51-phase-dictionary} and every integer
\(Q\ge2\), there is an integer \(q\) with
\(1\le q\le Q^N\) such that, for every \(t\in\R\),
\begin{equation}
 \|Z_{\bm y}(t+q)-Z_{\bm y}(t)\|_2
 \le\frac{2\pi}{Q}\|\bm a\|_2.
 \label{eq:tpc51-phase-alias-bound}
\end{equation}
Consequently the signed pair
\(\delta_t-\delta_{t+q}\) has condition number at least \(Q/\pi\).
\end{proposition}

\begin{proof}
Apply simultaneous Dirichlet approximation to
\(y_n/(2\pi)\): there is \(1\le q\le Q^N\) such that
\(\|qy_n/(2\pi)\|_{\R/\Z}\le Q^{-1}\) for every \(n\).
Hence

\[
 |\e^{iqy_n}-1|\le2\pi/Q.
\]

After factoring out the unit phases \(\e^{ity_n}\), summing the
squared coordinate bounds proves
\eqref{eq:tpc51-phase-alias-bound}.  The envelope of the signed pair is
\(2\|\bm a\|_2\).
\end{proof}

The alias \(q\) may be extremely large, where a fixed Schwartz Mellin
density is tiny.  Accordingly,
\cref{prop:tpc51-neighboring-atoms,prop:tpc51-simultaneous-alias}
disprove coefficient-uniform frame statements, but neither proposition
estimates the one canonical coefficient measure.  That fixed direction
is treated next.
